% Section 6: Discussion. New file.

Prices rose. Buyers paid them. Neither side reorganized. Five candidate mechanisms are consistent with this pattern; we evaluate each one against the data we have, rule one out, and leave the rest as hypotheses that the current data cannot separately identify.

\paragraph{1. Shock too small at the buyer's input bill.} The simplest explanation is that the policy-induced price gap between regulated and unregulated suppliers is too small to overcome any reasonable switching cost. The within-NACE-4d shock is large at the supplier level for cement, basic metals, and refining, but most buyers are not concentrated in these sectors; the distribution of the dollar shock at the buyer's total-cost level (Section~\ref{sec:domestic_within_intensive}) is heavily right-skewed, with the typical buyer-cell facing a shock that is an order of magnitude smaller than any plausible switching cost even at 2022 EUA prices. The heterogeneity tests in \S\ref{sec:domestic_within_intensive} and \S\ref{sec:domestic_across} address this directly by restricting to the top exposure quartile, where the shock is largest. The null persists. This rules out the strongest form of the shock-too-small story --- that substitution would appear in the cells where the shock bites hardest --- but leaves open a weaker form: even Q4 cells may face a shock too small relative to switching costs. We cannot separately identify switching costs from elasticities with the current data.

\paragraph{2. Strategic complementarities in pricing.} The relative-price signal that drives substitution is the gap between competing suppliers' prices, not the gap between their costs. If regulated suppliers raise prices in response to the carbon shock and their unregulated competitors raise prices in step --- because demand is partially common and pricing is strategically complementary --- the within-NACE-4d price gap will be smaller than the cost-shock gap. \citet{amitiitskhokionings2014} document this for Belgian exporters in the context of exchange-rate shocks. The mechanism predicts low cross-supplier price dispersion and partial pass-through at the buyer level, both of which are plausible in our data but neither of which we measure directly. We cannot rule this out.

\paragraph{3. Relational capital --- inconsistent with the data.} A widely-invoked explanation for low substitution in firm-to-firm networks is that buyer-supplier relationships carry value beyond price --- accumulated specific investments, trust, contract-specific information --- that makes substitution costly. If this were the binding friction in our setting, it would not bind on the new-relationship margin: a buyer forming a fresh supplier pair starts from a clean slate, with no prior investments to write off. Buyers forming new pairs in an ETS-treated NACE 4-digit sector after the MSR decision should therefore select systematically cleaner suppliers than buyers forming new pairs before. We test this directly, using the supplier's within-NACE 4-digit percentile rank of allowance shortage as the carbon-cost characteristic the buyer can observe at the time of pair formation (Appendix~\ref{app:omega_rank}). Restricting the sample to NACE 4-digit sectors where the within-sector supplier rank ordering is stable enough between $2016$ and $2022$ (Spearman~$\rho\geq 0.8$) that year-$t$ rank cleanly identifies supplier choice rather than within-firm rank evolution, the post-event event-study coefficients hover within $\pm 0.05$ of zero from rel-year~$+1$ onward, with $95\%$ confidence intervals that include zero (Figure~\ref{fig:omega_rank_stable_yeartrank}). The new-relationship margin --- the one margin where relational capital should bind most weakly --- shows no shift toward cleaner suppliers post-MSR. Read honestly given the thin EUTL universe ($\sim 28$ firms across $7$ stable sectors at the headline threshold) and imperfect pre-trends at $\tau=-2$ and $\tau=0$ (driven by the documented 2015--2016 B2B reporting-wave cohort effect, Appendix~\ref{app:omega_rank}), this is "no support for the relational-capital story" rather than a strict rejection. The data are inconsistent with relational capital being the dominant friction.

\paragraph{4. Partial pass-through.} The pass-through evidence in \S\ref{sec:passthrough} establishes that the EU ETS shock translates into the regulated PPI: the Känzig structural carbon-policy shock loads onto the Belgian PPI with peak coefficient $+4.08$ at horizon 12 months. The pass-through is not full, however --- the regulated-vs-unregulated PPI gap is $+0.094$ in Phase~II and $+0.113$ in Phase~IV, which corresponds to roughly 10 percentage points of differential price-level change, smaller than the underlying cost-shock differential. Partial pass-through dampens the relative-price signal that would induce buyer substitution. We have not estimated firm-level pass-through directly --- the customs panel records both quantities and values for cross-border transactions, but the B2B panel records euro amounts only --- so we cannot rule out a yet-smaller within-NACE-4d cross-supplier price wedge than the aggregate pass-through coefficients suggest. This mechanism remains consistent with the joint null.

\paragraph{5. Network-adjusted control-group contamination.} The unregulated suppliers in our within-NACE-4d sample are not actually clean: they purchase regulated inputs upstream and therefore inherit a share of the carbon shock through the production network. When emission intensity is network-adjusted (own emissions plus upstream emissions weighted by Leontief shares), the gap between the most-exposed regulated supplier and the ``unregulated alternative'' compresses. The relevant relative-price wedge for substitution may be much smaller than the direct emission-intensity gap suggests. Constructing the network-adjusted intensity requires a full Belgian input-output table at the firm level; we have the ingredients but have not implemented the test.

\paragraph{Reading the five together.} Of the five mechanisms, only relational capital is in tension with the data: the new-relationship margin --- where it should bind least --- shows no shift toward cleaner suppliers post-MSR. The other four --- shock-magnitude limits, strategic complementarities, partial pass-through, and network-adjusted contamination --- are all mechanisms that would compress the within-cell relative-price gap actually faced by Belgian buyers. They are not mutually exclusive, and in fact reinforce each other: a partially passed-through cost shock combined with strategic complementarities in pricing combined with a partially regulated control group leaves a small wedge for substitution to act on. Separating them empirically requires either direct firm-level price data (PRODCOM, deferred to a complementary paper with the co-author) or a richer model of pricing in buyer-supplier networks. We leave both to future work. What is settled is the headline: under the EU ETS as it currently operates in Belgium, the leakage that motivated the policy's defensive carve-outs is not measurable on any of the three margins where it could plausibly act.
